\SourceLeaf{230}%
\Article{5}
\indent The two solutions found in the preceding article evidently refer to opposite points on the spherical surface, or to opposite directions. This is natural, since the normal may be drawn toward either side of the curved surface. If we choose to distinguish the two regions adjoining the surface and call one exterior and the other interior, we can also assign each normal its appropriate solution by means of the theorem developed in Article 2 (VII), once a criterion has been established for distinguishing the two regions.

In the first method, such a criterion must be drawn from the sign of the value of \(W\). Generally speaking, the curved surface separates the parts of space in which \(W\) takes a positive value from those in which the value of \(W\) becomes negative. It readily follows from that theorem that if \(W\) takes a positive value toward the exterior side and we take the normal to be drawn outward, the first solution must be chosen. In each case it will be easy to decide whether the same rule for the sign of \(W\) holds throughout the whole surface or whether different rules hold for different parts. As long as the coefficients \(P,Q,R\) are finite and do not all three vanish at once, the law of continuity precludes any change.

If we follow the second method, we can conceive of two families of curves on the curved surface. In one, \(p\) varies while \(q\) is constant. In the other, \(q\) varies while \(p\) is constant. The relative position of these curves with respect to the exterior side must determine which solution to choose. Take the branch of a curve of the first family that leaves \(A\) as \(p\) increases, the branch of a curve of the second family that leaves \(A\) as \(q\) increases, and the normal drawn toward the exterior side. Whenever these three lines are arranged \textit{in the same way} as the respective axes \(x,y,z\) starting from the origin of the abscissas (for example, if in each set of three lines the first can be imagined as pointing left, the second right, and the third upward), the first solution must be chosen. Whenever the relative arrangement of the three lines is opposite to that of the axes \(x,y,z\), the second solution applies.

In the third method, we must consider whether giving \(z\) a positive increment while keeping \(x\) and \(y\) unchanged takes us toward the exterior or the interior. For the outward normal, the first solution applies in the former case, the second in the latter.
